\section{Numerical Implementation}\label{sec:numerical_impl}
% ============================================================
% ============================================================

\subsection{Variational method: algorithm}

The following steps summarize the complete numerical procedure implemented in QuTu:

\begin{enumerate}
  \item \textbf{Read physical parameters}: $\xe$, $\Vb$, reduced mass $\mu$
    (e.g., $\mu = 3m_H m_N/(3m_H+m_N)$ for NH$_3$), asymmetry parameter $\varepsilon$,
    and number of basis functions $N$.
  \item \textbf{Compute $\alpha_{\text{opt}}$} from \cref{sec:alpha_opt}:
    use the deep-well estimate $\alpha_{\text{opt}} = 2\sqrt{\mu b}$ as starting point,
    then refine numerically.
  \item \textbf{Symmetric case} ($\varepsilon = 0$): build $\mathbf{H}_{\text{even}}$
    (dimension $\lceil N/2\rceil$) and $\mathbf{H}_{\text{odd}}$ (dimension
    $\lfloor N/2\rfloor$) using \cref{eq:T_HO,eq:x2_HO,eq:x4_HO} restricted to
    even/odd quantum numbers; diagonalize both with a real symmetric eigensolver
    (e.g., LAPACK \texttt{DSYEV}); interleave eigenvalues as $E_0 < E_1 < E_2 < \cdots$
  \item \textbf{Asymmetric case} ($\varepsilon \neq 0$): build the full $N\times N$
    matrix $\mathbf{H}^{\text{asym}}$ from \cref{eq:H_asym} (adding
    $\varepsilon\mel{m}{\hat{x}}{n}$ from \cref{eq:x1_HO}); diagonalize with
    \texttt{DSYEV}.
  \item \textbf{Reconstruct wavefunctions}: from the eigenvectors
    $\{C_n^{(k)}\}$, evaluate $\Phi_k(x) = \sum_n C_n^{(k)}\varphi_n(x;\alpha)$
    on a real-space grid.
  \item \textbf{Construct localized wavepackets}: symmetric case via
    \cref{eq:Psi_pm}; asymmetric case via \cref{eq:psi_LR_asym,eq:mixing_angle}.
  \item \textbf{Compute observables}: survival probability from \cref{sec:Ps_general},
    $\langle\hat{x}\rangle(t)$ from \cref{sec:x_expect}, recurrence time
    $\trec = h/\Delta_0$, and additional observables as listed in
    \cref{sec:observables}.
\end{enumerate}

\subsection{Convergence criterion}

The energy $E_k^{(N)}$ computed with $N$ basis functions converges to the exact
value as $N\to\infty$. A practical convergence criterion is
\begin{equation}
  |E_k^{(N)} - E_k^{(N_{\max})}| < \delta_{\text{conv}}\,,
\end{equation}
where $N_{\max}$ is a large reference basis and $\delta_{\text{conv}}$ is a chosen
tolerance (e.g., $10^{-6}~E_h$). For the quartic double-well, the lowest levels
typically converge at $N \sim 20$--$50$.

\subsection{Unit conversion table}

The code works internally in atomic units. Key conversion factors:
\begin{center}
\renewcommand{\arraystretch}{1.3}
\begin{tabular}{lll}
  \toprule
  Quantity & Atomic unit & Practical unit \\
  \midrule
  Length   & $a_0 = 0.52918\;\text{\AA}$ & angstrom (\AA) \\
  Energy   & $E_h = 1\;\text{hartree}$   & $1\,E_h = 219474.63\;\text{cm}^{-1}$ \\
  Mass     & $m_e$                        & $1\;\text{u} = 1822.89\,m_e$ \\
  Time     & $\hb/E_h$                    & $1\;\hb/E_h = 24.19\;\text{as}$ \\
  \bottomrule
\end{tabular}
\end{center}

% ============================================================
% ============================================================

% ============================================================
% ============================================================
